\section{不放回抽样与 Hypergeometric 分布}
\label{sec:05-Probability/06-Common-Distributions/02}

你面前有一箱零件，共 \(100\) 个，质检员告诉你其中恰有 \(8\) 个次品——不多不少。你随手从箱子里抓出 \(10\) 个，不放回去，想根据这 \(10\) 个样品里次品的数量判断整批是否可接受。

如果次品比例真的是 \(8\%\)，每次抓取像是成功概率 \(0.08\) 的独立试验吗？显然不是。抓走一个次品后，剩下的次品比例从 \(8/100\) 变成了 \(7/99\)；抓走一个好零件，次品比例则升到 \(8/99\)。各次抽取结果互相依赖，普通 Binomial 模型在这里不对——它假设每次都是同一枚硬币抛掷，而你的硬币正在被自己一口一口吃掉。

这就是 Hypergeometric 分布的生长土壤：有限总体、固定组成、不放回抽样。它回答的问题是：已知罐子里白球黑球的精确数量，抓一把出来后，手里白球数量的分布。

\subsection{先数组合，再谈概率}

假设简单随机抽样，即所有大小为 \(n\) 的无序样本被抽到的可能性相同。全部可能的样本数：

\[
\binom{N}{n}.
\]

要样本中恰有 \(k\) 个成功对象，必须同时做两件事：从 \(K\) 个成功对象中选 \(k\) 个，从 \(N-K\) 个失败对象中选 \(n-k\) 个。因此计数为

\[
\binom{K}{k}\binom{N-K}{n-k},
\]

概率质量函数即

\[
P(X=k)=\frac{\binom{K}{k}\binom{N-K}{n-k}}{\binom{N}{n}}.
\]

这个公式看着干净，但写支撑时很多人栽跟头。你不能机械地写成 \(k=0,1,\ldots,n\)。\(k\) 的上限被两个条件同时约束：不能超过你抽的总数 \(n\)，也不能超过总体中实际存在的成功对象数 \(K\)。同理，下限也被约束：如果你抽了 \(n\) 个，而总体中失败对象一共只有 \(N-K\) 个，那你至少得抽出 \(n-(N-K)\) 个成功对象。支撑的准确写法是

\[
\max(0,\,n-(N-K))\le k\le\min(n,K).
\]

放进最开始那个零件的例子：\(N=100\)，\(K=8\)，\(n=10\)。\(k\) 的上限是 \(\min(10,8)=8\)，下限是 \(\max(0,10-92)=0\)。支撑确实是 \(0\) 到 \(8\)——你不可能抽出 \(9\) 或 \(10\) 个次品，因为箱子里一共就 \(8\) 个。

如果把各 \(k\) 对应的概率逐个算一遍，分布的形状就出来了。\(k=0\) 时——一个次品都没抓到——概率是 \(\binom{8}{0}\binom{92}{10}/\binom{100}{10}\approx0.416\)。\(k=1\) 约为 \(0.393\)。\(k=2\) 约为 \(0.147\)。概率随 \(k\) 增加迅速衰减：抓到 \(4\) 个次品的概率已经降到约 \(0.006\)，\(k=8\) 更是远低于千分之一。大部分抽样结果集中在 \(0\) 到 \(2\) 个次品之间，这和次品率 \(8\%\) 的印象是协调的——抽 \(10\) 个，期望抓到 \(0.8\) 个，但分布有明显的右尾巴。

\subsection{均值为啥和 Binomial 一样}

尽管各次抽取不独立，每个特定位置的边缘行为仍然简单。令 \(I_i\) 指示第 \(i\) 个抽出的是否为成功对象。在简单随机抽样下，每个对象被抽到任一特定位置的几率相同，由对称性：

\[
P(I_i=1)=\frac{K}{N},\qquad i=1,\ldots,n.
\]

记 \(p=K/N\)。由期望的线性性——这条永远不需要独立：

\[
\E[X]=\E\left[\sum_{i=1}^{n}I_i\right]=\sum_{i=1}^{n}\frac{K}{N}=np.
\]

期望和 Binomial\((n,p)\) 一模一样。这是因为``平均来看''每次抽取的成功机会并没有被不放回改变——不是独立不独立决定的，是边缘概率决定的。期望的线性性是概率论中最忠实的工具：只要你求和的对象定义清楚了，依赖关系一概不管，期望照加不误。

\subsection{方差：依赖的代价}

方差就不一样了。因为不放回引入了负相关：抽到一个成功对象，说明``好运气用掉了一点''，后续成功的概率微降。

对 \(i\neq j\)，

\[
\Cov(I_i,I_j)=P(I_i=I_j=1)-p^2
=\frac{K(K-1)}{N(N-1)}-p^2.
\]

这个协方差不等于零。代入 \(p=K/N\) 并整理：

\[
\begin{aligned}
\Cov(I_i,I_j)
&=\frac{K(K-1)}{N(N-1)}-\frac{K^2}{N^2}\\[2pt]
&=\frac{K}{N}\left(\frac{K-1}{N-1}-\frac{K}{N}\right)\\[2pt]
&=\frac{K}{N}\cdot\frac{N(K-1)-K(N-1)}{N(N-1)}\\[2pt]
&=\frac{K}{N}\cdot\frac{-(N-K)}{N(N-1)}\\[2pt]
&=-\frac{p(1-p)}{N-1}.
\end{aligned}
\]

负号是依赖方向的定量表达。协方差对所有 \(i\neq j\) 都是同一个常数——任意两次抽取之间的相互牵制强度相同。现在把方差展开：

\[
\begin{aligned}
\Var(X)&=\sum_{i=1}^{n}\Var(I_i)+\sum_{i\neq j}\Cov(I_i,I_j)\\[2pt]
&=np(1-p)+n(n-1)\left(-\frac{p(1-p)}{N-1}\right)\\[2pt]
&=np(1-p)\left(1-\frac{n-1}{N-1}\right)\\[2pt]
&=n\,\frac{K}{N}\left(1-\frac{K}{N}\right)\frac{N-n}{N-1}.
\end{aligned}
\]

最右边的因子

\[
\frac{N-n}{N-1}
\]

就是有限总体修正（finite population correction）。它的值在 \(0\) 到 \(1\) 之间，衡量不放回相对于放回把方差压掉了多少。

两个极端验证你的直觉：若 \(n=1\)——只抽一个，无所谓放回不放回——修正因子为 \((N-1)/(N-1)=1\)，方差退化为 \(p(1-p)\)，与 Bernoulli 一致。若 \(n=N\)——把整个总体提走——成功数必为 \(K\)，方差应为零，而此时修正因子恰为 \((N-N)/(N-1)=0\)。公式两端都说得通。

回到零件例子感受一下数字。Binomial\((10,0.08)\) 的方差是 \(10\times0.08\times0.92=0.736\)。Hypergeometric 的方差是 \(0.736\times(100-10)/(100-1)\approx0.736\times0.909=0.669\)。不放回把方差缩小了约 \(9\%\)——不算剧烈，但如果你抽的不是 \(10\) 个而是 \(50\) 个，Binomial 方差 \(50\times0.08\times0.92=3.68\)，而 Hypergeometric 方差要乘上 \((100-50)/(100-1)\approx0.505\)，只有约 \(1.86\)，差了一倍。抽样越多，两种模型的分歧越大。

\subsection{Binomial 近似：\"足够大\"有多\"大\"}

当总体很大而你只摸了很小一部分时，不放回对组成的扰动微乎其微。此时

\[
\frac{N-n}{N-1}\approx 1,
\]

Hypergeometric 可以用

\[
\mathrm{Binomial}\!\left(n,\frac{K}{N}\right)
\]

近似。常见经验说法是抽样比例不超过 \(5\%\) 或 \(10\%\)。但这个数字没有普适的数学效力——它只在方差近似误差在你容忍范围内时才成立，而你的容忍范围取决于你接下来要用这个分布做什么决策。

判断近似质量可以直接看有限总体修正因子。举几个数值对照：抽样比例 \(n/N=5\%\)，修正因子约 \(0.95\)，方差低估约 \(5\%\)。\(n/N=20\%\)，修正因子约 \(0.80\)，方差低估约 \(20\%\)。\(n/N=50\%\)，修正因子约 \(0.505\)，方差低估约一半——此时把 Hypergeometric 当 Binomial 用，你会以为自己有原来两倍的不确定性。这个误差意味着置信区间被不必要地拉宽，假设检验的功效被低估。

还有一个更隐蔽的问题：即使方差近似尚可，Hypergeometric 的尾部行为和 Binomial 仍有差别。因为总体中成功对象数量有硬上限，Hypergeometric 的右尾在接近 \(K\) 时会被截断，而 Binomial 的正态近似则允许理论上任意大的值。做极端尾部概率计算时，Binomial 近似可能给出物理上不可能的结论。

\subsection{区别不在\"是否放回\"，在\"随机性的来源\"}

放回是 Binomial 模型最自然的物理实现；不放回是 Hypergeometric 最自然的物理实现。但二选一不应只看物理操作。

更根本的诊断是：成功对象的数量在抽样过程中是固定的还是随机的？

\begin{itemize}
\tightlist
\item
  Hypergeometric 的假设：总体中成功对象的个数已经确定，只是你不知道抽样会抽出哪些。你的不确定性来自抽样本身的随机性。
\item
  Binomial 的假设：每个对象有一个固定的成功概率，各次试验独立。你的不确定性来自每次试验的随机结果，并且成功对象的总数本身也是随机的。
\end{itemize}

用制造场景区分：如果一条产线每分钟稳定地以 \(8\%\) 的缺陷率产出零件，你抽 \(10\) 个，每个零件独立地有 \(0.08\) 概率为次品，适合用 Binomial。如果一批已经生产完毕的 \(100\) 个零件中，质检员盘点后发现有 \(8\) 个次品隐藏其中，你要抽 \(10\) 个——适合用 Hypergeometric。

数据表面相同：都是 \(100\) 个，\(8\%\) 的次品率，抽 \(10\) 个。但数据生成的物理机制不同，对应的概率模型就不同。分清楚随机性到底在哪个环节进入，比记住``放回用 Binomial、不放回用 Hypergeometric''这个口诀可靠得多。

换个场景让二者的分野更尖锐。选民调查：全市有 \(50\) 万注册选民，你随机拨打 \(1000\) 个电话，问支持谁。选民总数固定，你的抽样不放回，但因为 \(n/N=1000/500000=0.2\%\)，有限总体修正几乎等于 \(1\)，Binomial 近似精准。但如果你调查的是一个 \(30\) 人的班级，对 \(25\) 人做了访谈，此时 \(n/N=83\%\)，修正因子约 \(0.17\)——必须老老实实用 Hypergeometric，不然方差会高估将近六倍。

\subsection{回到那个零件例子，算个数}

\(N=100\)，\(K=8\)，\(n=10\)。恰发现 \(2\) 个次品的概率：

\[
P(X=2)=\frac{\binom{8}{2}\binom{92}{8}}{\binom{100}{10}}.
\]

分子第一部分 \(\binom{8}{2}=28\)：从 \(8\) 个次品中挑出你要抓到的那 \(2\) 个。第二部分 \(\binom{92}{8}\)：从 \(92\) 个好零件选 \(8\) 个。分母 \(\binom{100}{10}\) 是总组合数。

如果错用 Binomial\((10,0.08)\)，算出来的是 \(P_{\text{Bin}}(X=2)=\binom{10}{2}(0.08)^2(0.92)^8\approx0.148\)，而 Hypergeometric 精确值也在 \(0.147\) 左右——两个数字很接近，抽样比例只有 \(10\%\)，修正因子 \((100-10)/(100-1)=90/99\approx0.909\)，还不够剧烈。但如果你抽的是 \(50\) 个，差异就会非常显著。关键是知道哪个模型对应你面对的真实生成机制，而不是盲目套用更简单的那一个。

\subsection{从计数到检验：Fisher 精确检验的土壤}

Hypergeometric 分布有一个直通统计推断的入口。回到那个典型场景：把 \(N\) 个对象按两个维度交叉分类，比如治疗组/对照组和治愈/未治愈。固定所有边际——两组各多少人、总共治愈多少——问在无关联的零假设下，治疗组的治愈人数服从什么分布。

答案正是 Hypergeometric。把治愈的人视为``成功对象''，从全体中不放回地分配到治疗组。Fisher 精确检验的 \(p\) 值就是这个 Hypergeometric 分布的尾部概率。它和本节 PMF 的关系只有一步之遥：检验只是把计数的分布反过来用于判断关联。

\subsection{相关但不相同的帽子}

不放回抽样还连接着几个近邻概念。Coupon collector 问题研究的是``每个不同的对象至少要抽到一次需要多少次''——它关心覆盖而非计数。多变量 Hypergeometric 把成功/失败的二分推广到多类别：罐子里有红蓝绿三种球，一次抓一把，问各种颜色分别抓到多少；其 PMF 分子变成多个二项式系数的乘积。中心极限定理下，大样本的 Hypergeometric 计数也可以用 Gaussian 近似，但有限总体修正仍要保留——近似的是``放回式''Gaussian 乘上修正因子，而非直接套用 Binomial 的正态近似。

Hypergeometric 给你的核心提醒是：当总体是固定有限集合时，每次观察会改变剩余部分的信息状态。这个效应在总量小或抽样比例高时不可忽视——它不是一个可以等学了渐近理论再回来打补丁的细节。下一次你面对``放回还是不放回''的选择时，先问自己：随机性到底在哪里——是在抽样的手上，还是在对象的生成机制里？
